\documentclass[11pt]{amsart}

\usepackage[T1]{fontenc}
\usepackage{lmodern}
\usepackage{microtype}
\usepackage{amsmath,amssymb}
\usepackage{booktabs}
\usepackage[hidelinks]{hyperref}

\hypersetup{
  pdftitle={A Counterexample to a Cartesian-Product Conjecture for Node Resistance Curvature},
  pdfsubject={Node resistance curvature in Cartesian graph products},
  pdfkeywords={effective resistance, node resistance curvature, Cartesian product, spanning tree}
}

\newtheorem{theorem}{Theorem}
\newtheorem*{conjecture}{Conjecture 4}
\theoremstyle{remark}
\newtheorem*{remark}{Remark}

\title[Counterexample for node resistance curvature]
{A Counterexample to a Cartesian-Product Conjecture\\
for Node Resistance Curvature}

\date{}

\subjclass[2020]{05C50, 05C76}
\keywords{effective resistance, node resistance curvature, Cartesian product,
spanning tree}

\begin{document}

\begin{abstract}
We disprove Conjecture~4 of Dawkins et al. as stated in
arXiv:2403.01037v1. Two rooted trees on $5$ and $8$ vertices have root
curvature zero, whereas the corresponding vertex in their Cartesian
product has node resistance curvature
$32552200/4351396379>0$. The verification uses exact integer
arithmetic and the matrix-tree theorem.
\end{abstract}

\maketitle

\section{The counterexample}

Let $G=(V,E)$ be a finite connected simple graph with unit edge
conductances, let $L_G$ be its Laplacian, and let $L_G^+$ be the
Moore--Penrose pseudoinverse of $L_G$. For $u,v\in V$, write
\[
\omega_G(u,v)
=(\mathbf e_u-\mathbf e_v)^{\mathsf T}
L_G^+(\mathbf e_u-\mathbf e_v)
\]
for the effective resistance. The node resistance curvature introduced
by Devriendt and Lambiotte~\cite{DL} is
\[
p_v(G)=1-\frac12\sum_{u\sim v}\omega_G(u,v).
\]
We write $G_1\square G_2$ for the Cartesian product.

In arXiv:2403.01037v1, Dawkins et al.~\cite{DGKLQRS} state the
following.

\begin{conjecture}
For graphs $G_1,G_2$ and vertices $i\in V(G_1)$ and $j\in V(G_2)$, if
\[
p_i(G_1)\le 0
\qquad\text{and}\qquad
p_j(G_2)\le 0,
\]
then
\[
p_{(i,j)}(G_1\square G_2)<0.
\]
\end{conjecture}

\begin{theorem}
Let $G_1$ have vertex set $\{0,1,2,3,4\}$ and edge set
\[
E(G_1)=\{01,02,13,34\},
\]
and let $G_2$ have vertex set $\{0,1,\ldots,7\}$ and edge set
\[
E(G_2)=\{01,02,03,04,15,56,57\},
\]
where $ab$ denotes the edge $\{a,b\}$. At the marked vertices $i=j=1$,
\[
p_i(G_1)=p_j(G_2)=0,
\]
but
\[
p_{(i,j)}(G_1\square G_2)
=\frac{32552200}{4351396379}>0.
\]
Consequently, Conjecture~4 is false as stated.
\end{theorem}

\begin{proof}
Both $G_1$ and $G_2$ are trees, and the marked vertex has degree $2$ in
each. Every unit edge of a tree has effective resistance $1$, so
\[
p_i(G_1)=p_j(G_2)=1-\frac22=0.
\]

Set $H=G_1\square G_2$ and $x=(1,1)$. The neighbors of $x$ are
\[
(0,1),\qquad (3,1),\qquad (1,0),\qquad (1,5).
\]
For a connected graph $K$, let $\tau(K)$ denote its number of spanning
trees. If $e=uv$ is a unit edge, then
\[
\omega_K(u,v)=\Pr(e\in\mathbf T)=\frac{\tau(K/e)}{\tau(K)},
\]
where $\mathbf T$ is a uniformly random spanning tree and $K/e$ is
obtained by contracting $e$, with parallel edges retained with
multiplicity. The first equality is the spanning-tree inclusion
identity~\cite[Theorem~2]{DL}, and the second follows from the standard
contraction bijection.

Applying the matrix-tree theorem to the reduced Laplacians gives
\[
\tau(H)=723{,}963{,}572{,}556{,}125
\]
and the following four cofactors.
\begin{center}
\small
\renewcommand{\arraystretch}{1.25}
\begin{tabular}{@{}lr@{}}
\toprule
contracted edge $e$ & $\tau(H/e)$ \\
\midrule
$\{x,(0,1)\}$ & $359{,}244{,}778{,}688{,}125$ \\
$\{x,(3,1)\}$ & $359{,}244{,}778{,}688{,}125$ \\
$\{x,(1,0)\}$ & $354{,}728{,}028{,}191{,}625$ \\
$\{x,(1,5)\}$ & $363{,}877{,}814{,}994{,}375$ \\
\bottomrule
\end{tabular}
\end{center}
Therefore
\[
\begin{aligned}
p_x(H)
&=\frac{2\tau(H)-\sum_{e\ni x}\tau(H/e)}{2\tau(H)}\\
&=\frac{10{,}831{,}744{,}550{,}000}
        {1{,}447{,}927{,}145{,}112{,}250}\\
&=\frac{32552200}{4351396379}>0.
\end{aligned}
\]
This contradicts the asserted strict negativity.
\end{proof}

\begin{remark}
The counterexample uses equality in both hypotheses. It therefore
refutes Conjecture~4 as stated, but it does not address the variant with
$p_i(G_1)<0$ and $p_j(G_2)<0$. No minimality claim is made.
\end{remark}

\begin{thebibliography}{9}

\bibitem{DGKLQRS}
A.~Dawkins, V.~Gupta, M.~Kempton, W.~Linz, J.~Quail, D.~H.~Richman,
and Z.~Stier,
\newblock \emph{Node resistance curvature in Cartesian graph products},
\newblock Journal of Combinatorics \textbf{16} (2025), no.~4, 575--589.
\newblock
\href{https://doi.org/10.4310/JOC.250923003413}
{doi:10.4310/JOC.250923003413};
\href{https://arxiv.org/abs/2403.01037v1}{arXiv:2403.01037v1}.

\bibitem{DL}
K.~Devriendt and R.~Lambiotte,
\newblock \emph{Discrete curvature on graphs from the effective resistance},
\newblock Journal of Physics: Complexity \textbf{3} (2022), no.~2,
025008.
\newblock
\href{https://doi.org/10.1088/2632-072X/ac730d}
{doi:10.1088/2632-072X/ac730d}.

\end{thebibliography}

\end{document}